\documentclass[UTF8,10pt,landscape]{ctexart}

% Build with: xelatex cheatsheet-template.tex
% Or:         lualatex cheatsheet-template.tex

\usepackage[
  a4paper,
  landscape,
  left=6mm,
  right=6mm,
  top=6mm,
  bottom=6mm
]{geometry}
\usepackage{multicol}
\usepackage{amsmath,amssymb,mathtools}
\usepackage{booktabs,tabularx,array}
\usepackage[table]{xcolor}
\usepackage{graphicx}
\usepackage{listings}
\usepackage[hidelinks]{hyperref}

\definecolor{accent}{HTML}{0F766E}
\definecolor{accentlight}{HTML}{E6FFFA}
\definecolor{ink}{HTML}{111827}
\definecolor{muted}{HTML}{6B7280}
\definecolor{rulegray}{HTML}{D1D5DB}
\definecolor{codebg}{HTML}{F8FAFC}

\newcommand{\cheatbodyfont}{\fontsize{8.6pt}{9.2pt}\selectfont}
\newcommand{\cheatsectionfont}{\fontsize{9.6pt}{10.1pt}\selectfont}
\newcommand{\cheatsubsectionfont}{\fontsize{8.8pt}{9.4pt}\selectfont}
\newcommand{\cheattinyfont}{\fontsize{7.7pt}{8.2pt}\selectfont}

\pagestyle{empty}

\setlength{\parindent}{0pt}
\setlength{\parskip}{0.4pt}
\setlength{\columnsep}{3.2mm}
\setlength{\columnseprule}{0.15pt}
\setlength{\premulticols}{0pt}
\setlength{\postmulticols}{0pt}
\setlength{\multicolsep}{0pt}
\setlength{\tabcolsep}{2pt}
\setlength{\fboxsep}{1pt}
\renewcommand{\arraystretch}{0.82}

\makeatletter
\def\@listi{\leftmargin\leftmargini
  \topsep 0.4pt
  \parsep 0pt
  \partopsep 0pt
  \itemsep 0pt}
\let\@listI\@listi
\@listi
\makeatother

\newcommand{\cheatsection}[1]{%
  \vspace{3pt}%
  {\color{accent}\normalfont\bfseries\cheatsectionfont #1\par}%
  \vspace{-2.5pt}%
  {\color{accent}\rule{\linewidth}{0.3pt}\par}%
  \vspace{1pt}%
}
\newcommand{\cheatsubsection}[1]{%
  \vspace{1.6pt}%
  {\color{ink}\normalfont\bfseries\cheatsubsectionfont #1\par}%
  \vspace{0.2pt}%
}

\newcommand{\zhhint}[1]{{\color{muted}\cheattinyfont（#1）}}
\newcommand{\term}[2][]{%
  \textbf{\color{accent}#2}%
  \if\relax\detokenize{#1}\relax
  \else\zhhint{#1}%
  \fi
}
\newcommand{\key}[1]{\fcolorbox{rulegray}{accentlight}{\cheattinyfont\texttt{#1}}}
\newcommand{\code}[1]{\texttt{\cheatbodyfont #1}}
\newcommand{\note}[1]{%
  \begingroup
  \setlength{\fboxsep}{2pt}%
  \colorbox{accentlight}{%
    \parbox{\dimexpr\linewidth-2\fboxsep\relax}{\cheattinyfont #1}%
  }%
  \endgroup
}
\newenvironment{stemlist}{%
  \begingroup
  \cheattinyfont
  \setlength{\parskip}{0.6pt}%
}{%
  \par
  \endgroup
}
\newcommand{\stemrow}[3][]{\par\noindent\term[#1]{#2}: #3\par}

\newcolumntype{Y}{>{\raggedright\arraybackslash}X}
\newcolumntype{T}{>{\raggedright\arraybackslash}p{0.34\linewidth}}
\newcolumntype{S}{>{\raggedright\arraybackslash}p{0.18\linewidth}}

\lstdefinestyle{cheatcode}{
  basicstyle=\ttfamily\cheattinyfont,
  backgroundcolor=\color{codebg},
  columns=fullflexible,
  keepspaces=true,
  breaklines=true,
  frame=single,
  rulecolor=\color{rulegray},
  framerule=0.2pt,
  xleftmargin=1pt,
  xrightmargin=1pt,
  aboveskip=1pt,
  belowskip=1pt
}
\lstset{style=cheatcode}

\begin{document}
\cheatbodyfont
\setlength{\abovedisplayskip}{1pt}
\setlength{\belowdisplayskip}{1pt}
\setlength{\abovedisplayshortskip}{0pt}
\setlength{\belowdisplayshortskip}{0pt}

\begin{multicols*}{4}
\cheatsection{0 Quick Index / 快速索引}
\begin{stemlist}
\term[总览/中断/缓存]{Ch1}: CPU registers, interrupts, memory hierarchy, cache, DMA.\\
\term[OS目标/接口/模式]{Ch2}: convenience, efficiency, evolution; ISA/ABI/API; system call; user/kernel mode.\\
\term[进程/状态/切换]{Ch3}: process image, PCB, five/seven-state model, suspend, trap, process switch.\\
\term[线程模型]{Ch4}: thread vs process, ULT/KLT, jacketing, shared resources.\\
\term[并发同步]{Ch5}: race condition, critical section, mutex, semaphore, producer/consumer.\\
\term[死锁处理]{Ch6}: four conditions, prevention, avoidance, detection, banker, safe state.\\
\term[内存管理]{Ch7}: relocation, partitioning, fragmentation, buddy, paging, address translation.\\
\term[虚拟内存]{Ch8}: page fault, resident set, TLB, OPT/LRU/FIFO/Clock, working set, thrashing.\\
\term[单处理器调度]{Ch9}: long/medium/short scheduler, FCFS, RR, SPN/SJF, SRT, HRRN, feedback.\\
\term[多核/实时]{Ch10}: load sharing, gang scheduling, dedicated/dynamic scheduling, hard/soft real-time.\\
\term[文件管理]{Ch12}: file organization, directory, B-tree, blocking, contiguous/chained/indexed allocation.
\end{stemlist}
\note{考试查找路线：先按章定位，再扫绿色关键词后的中文括注；计算题优先看 Ch7 地址转换、Ch8 置换/TLB、Ch9 调度甘特图、Ch12 B-tree/分配。}

\cheatsection{1 Computer System Overview}
\cheatsubsection{Main Elements and Registers}
\begin{tabularx}{\linewidth}{@{}TY@{}}
\toprule
\term[处理器/CPU]{Processor} & executes instructions and controls system operation.\\
\term[主存/当前程序数据]{Main memory} & stores currently used programs and data.\\
\term[输入输出模块]{I/O modules} & move data between the computer and external devices.\\
\term[数据/地址/控制线]{System bus} & carries data, address, and control signals among modules.\\
\bottomrule
\end{tabularx}
\term[用户可见寄存器]{User-visible registers}: referenced by machine instructions to reduce memory access.\\
\term[控制/状态寄存器]{Control/status registers}: PC, IR, MAR, MBR, PSW; used by CPU and OS to control execution.

\cheatsubsection{Instruction Actions}
\begin{itemize}
  \item \term[访存指令]{Processor-memory}: load/store between CPU and memory.
  \item \term[I/O传送]{Processor-I/O}: transfer data between CPU and I/O modules.
  \item \term[算术/逻辑]{Data processing}: arithmetic or logical operations.
  \item \term[改变执行流]{Control}: branch, jump, call, return; changes execution order.
\end{itemize}

\cheatsection{Interrupts}
\cheatsubsection{Purpose and Types}
\term[中断/提高CPU利用率]{Interrupt}: a mechanism by which I/O, timer, program error, or hardware failure interrupts normal processor execution. Main purpose: improve processor utilization.
\begin{tabularx}{\linewidth}{@{}TY@{}}
\toprule
\term[程序异常]{Program} & divide by zero, overflow, illegal instruction.\\
\term[时钟中断]{Timer} & lets the OS regain control periodically.\\
\term[设备完成/就绪]{I/O} & device reports completion or readiness.\\
\term[硬件故障]{Hardware failure} & power failure, memory parity error.\\
\bottomrule
\end{tabularx}

\cheatsubsection{Processing Sequence}
\begin{enumerate}
  \item Device issues interrupt; CPU finishes current instruction.
  \item CPU acknowledges, saves PC and PSW, and loads ISR address.
  \item ISR saves remaining state, handles the event, restores state.
  \item Old PC and PSW are restored; interrupted program resumes.
\end{enumerate}
\note{PC tells where to resume; PSW records processor state. Polling repeatedly tests device flags; interrupts let the device notify the CPU.}

\cheatsubsection{Multiple Interrupts}
\term[顺序处理]{Sequential}: disable interrupts during ISR; new interrupts stay pending. Simple but may delay urgent events.\\
\term[嵌套/优先级]{Nested}: assign priorities; a higher-priority interrupt may interrupt a lower-priority ISR.

\cheatsection{Memory, Cache, I/O}
\cheatsubsection{Memory Hierarchy}
Higher levels are \term[更快]{faster}, \term[更小]{smaller}, and \term[每位更贵]{more expensive per bit}; lower levels are larger and slower. Distinguish levels by capacity, access time, cost per bit, and frequency of access.
\[
T_{\mathrm{avg}}=\sum_i\left(\prod_{j<i}(1-H_j)\right)T_i
\]
\[
C_{\mathrm{avg}}=\sum C_iS_i/\sum S_i
\]

\cheatsubsection{Cache and Locality}
\term[高速缓存]{Cache}: small fast memory between processor and main memory; stores recently or frequently used blocks to reduce average access time.\\
\term[时间局部性/最近还会用]{Temporal locality}: a referenced item is likely to be referenced again soon; keep recent items in fast memory.\\
\term[空间局部性/附近也会用]{Spatial locality}: nearby addresses are likely to be referenced soon; fetch a block, not only one word.

\cheatsubsection{I/O and DMA}
\term[程序控制I/O/轮询]{Programmed I/O}: CPU waits by polling flags such as FGI/FGO.\\
\term[中断驱动I/O]{Interrupt-driven I/O}: IEN enables the device to interrupt when ready.\\
\term[直接存储器访问]{DMA}: transfers data directly between I/O device and memory; DMA gets high memory priority because device delay may cause data loss, while CPU delay usually only slows execution.

\cheatsection{Chapter 1 Exam Questions}
\cheatsubsection{Review 1.1--1.10}
\term[四大部件]{Q}: List the four main computer elements.\\
\term[CPU/内存/I/O/总线]{A}: Processor, main memory, I/O modules, system bus.\\
\term[寄存器分类]{Q}: Define the two register categories.\\
\term[用户可见+控制状态]{A}: User-visible registers; control and status registers.\\
\term[中断定义]{Q}: What is an interrupt?\\
\term[事件打断CPU转ISR]{A}: A mechanism that lets another module interrupt normal CPU execution so an ISR can handle the event.\\
\term[多中断处理]{Q}: How are multiple interrupts handled?\\
\term[顺序或嵌套]{A}: Sequential processing disables interrupts during ISR; nested processing uses priorities.\\
\term[存储层次指标]{Q}: What distinguishes memory hierarchy levels?\\
\term[容量/时间/成本/频率]{A}: Capacity, access time, cost per bit, and access frequency.\\
\term[多处理器vs多核]{Q}: Multiprocessor vs multicore?\\
\term[多个CPU vs 单芯片多核]{A}: Multiprocessor has multiple processors; multicore has multiple cores on one chip.

\cheatsubsection{Problem Stems and Answers}
\begin{stemlist}
\stemrow[I/O指令结果]{1.1 I/O program}{Load dev5, add M[940], store dev6: dev6 receives \code{0005}.}
\stemrow[取指/存储器寄存器]{1.2 MAR/MBR}{Expand Fig. 1.4 execution: fetch via MAR/MBR; final M[941] = \code{0005}.}
\stemrow[地址位数算容量]{1.3 32-bit instr.}{24-bit address field gives 16 MB; PC = 24 bits, IR = 32 bits.}
\stemrow[地址空间换算]{1.4 16-bit addr.}{16-bit memory: 128 KiB; 8-bit memory: 64 KiB; 8-bit port number: 256 ports.}
\stemrow[总线吞吐量]{1.5 bus rate}{8 MHz, 16-bit bus, 4 clocks/bus cycle gives 4 MB/s; wider bus or doubled clock gives 8 MB/s.}
\stemrow[轮询改中断]{1.6 Teletype I/O}{FGI/FGO polling wastes CPU; IEN enables interrupt-driven I/O.}
\stemrow[DMA优先权]{1.7 DMA priority}{DMA gets memory priority because device delay may lose data; CPU delay mostly slows execution.}
\stemrow[DMA占用比例]{1.8 DMA slowdown}{9600 bps DMA slows 1 MIPS CPU by about 0.12\%.}
\stemrow[I/O速率比较]{1.9 I/O rate}{Programmed I/O = 25,000 words/s; DMA = \(2.15\times10^6\) words/s.}
\stemrow[空间/时间局部性]{1.10 locality code}{Spatial: sequential \code{a[i]}; temporal: repeated same \code{a[i]} in inner loop.}
\stemrow[多级存储公式]{1.11 n-level memory}{Use \(T_{\mathrm{avg}}=\sum_i(\prod_{j<i}(1-H_j))T_i\), \(C_{\mathrm{avg}}=\sum C_iS_i/\sum S_i\).}
\stemrow[成本和命中率]{1.12 cost/hit}{Main memory \(\approx \$83.89\); cache-tech \(\approx \$838.86\); \(H\approx99.17\%\).}
\stemrow[平均访问时间]{1.13 cache/disk}{Average access time \(\approx480{,}026\) ns.}
\stemrow[PC不能省]{1.14 stack PC}{No; the PC cannot generally be eliminated by using the stack top.}
\end{stemlist}

\cheatsection{2.1 OS Objectives and Functions}
\cheatsubsection{Definition}
\term[操作系统]{Operating system}: a program that controls application execution and acts as an interface between applications and computer hardware.

\cheatsubsection{Three Objectives}
\begin{tabularx}{\linewidth}{@{}TY@{}}
\toprule
\term[方便使用]{Convenience} & makes the computer easier to use.\\
\term[资源效率]{Efficiency} & uses CPU, memory, I/O, files, storage, and network resources efficiently.\\
\term[可演化/可维护]{Evolution} & permits new functions, tests, fixes, and hardware support without disrupting service.\\
\bottomrule
\end{tabularx}

\cheatsubsection{Two Main Roles}
\begin{itemize}
  \item \term[用户接口]{User/computer interface}: hides hardware details and provides usable services.
  \item \term[资源管理器]{Resource manager}: allocates processor, memory, I/O devices, files, storage, and network resources.
\end{itemize}

\cheatsubsection{OS Services}
\begin{tabularx}{\linewidth}{@{}TY@{}}
\toprule
\term[程序开发服务]{Program development} & editors, debuggers, compilers, libraries.\\
\term[程序执行服务]{Program execution} & loads code/data, initializes files, I/O, and resources.\\
\term[I/O统一访问]{I/O access} & uniform device interface, e.g. \code{read()}, \code{write()}.\\
\term[文件访问]{File access} & file organization, permissions, protection.\\
\term[系统访问/权限]{System access} & login, authorization, resource contention control.\\
\term[错误处理]{Error handling} & detect hardware/software errors; terminate, retry, or report.\\
\term[记账/性能统计]{Accounting} & collect usage/performance data; support billing and tuning.\\
\bottomrule
\end{tabularx}

\cheatsubsection{Interfaces: ISA vs ABI vs API}
\begin{tabularx}{\linewidth}{@{}SYY@{}}
\toprule
\term[指令集接口]{ISA} & CPU instruction set & hardware/software boundary; user ISA vs system ISA.\\
\term[二进制接口]{ABI} & binary compatibility & system-call interface, registers, calling convention, binary format.\\
\term[源代码接口]{API} & source-level calls & library functions used by programmers, e.g. \code{open()}, \code{malloc()}.\\
\bottomrule
\end{tabularx}
\note{Memory hook: API = coding; ABI = compiled binary execution; ISA = CPU instructions.}

\cheatsection{Kernel, Control, Evolution}
\cheatsubsection{Kernel / Nucleus}
\term[内核/常驻核心]{Kernel}: central OS part, normally resident in main memory, containing the most frequently used OS functions: process, memory, I/O, file, interrupt, and system-call handling.

\cheatsubsection{How the OS Controls the CPU}
The OS is also software. It gains control, schedules and allocates resources, gives control to user programs, then regains control through interrupts, system calls, or exceptions.

\cheatsubsection{Why OSs Must Evolve}
\begin{itemize}
  \item \term[硬件升级]{Hardware upgrades}: new processors/devices require new support.
  \item \term[新增服务]{New services}: users and administrators need new tools and features.
  \item \term[修复错误]{Fixes}: bugs must be repaired without destabilizing old services.
\end{itemize}

\cheatsection{System Calls and Modes}
\cheatsubsection{System Calls}
\term[系统调用/进内核]{System call}: controlled entry point for a user program to request OS services such as file I/O, process creation, memory allocation, communication, and protection.

\cheatsubsection{Dual-Mode Operation}
\begin{tabularx}{\linewidth}{@{}TY@{}}
\toprule
\term[用户态/低权限]{User mode} & low privilege; applications cannot directly execute privileged instructions.\\
\term[内核态/高权限]{Kernel mode} & high privilege; OS kernel performs protected operations.\\
\bottomrule
\end{tabularx}
\note{Flow: a user program requests a system call/trap; the kernel performs the service and returns to user mode.}

\cheatsection{PDF Stems and Answers}
\begin{stemlist}
\stemrow[OS三目标]{R2.1}{Stem: three OS-design objectives. A: convenience, efficiency, evolution.}
\stemrow[内核定义]{R2.2}{Stem: kernel of an OS. A: memory-resident core with frequent OS functions.}
\stemrow[系统调用+双模式]{P2.4}{Stem: system-call purpose and dual-mode relation. A: safe service request; user mode traps to kernel mode, then returns.}
\stemrow[页框检查/内存压力]{P2.5}{Stem: OS/390 SRM checks untouched page frames. A: find inactive pages, guide replacement/reallocation, detect memory pressure or thrashing.}
\stemrow[存储管理职责]{R2.6}{Stem: five storage duties. A: isolation, allocation, modular support, protection/access control, long-term storage.}
\end{stemlist}

\cheatsection{Chapter 2 Fast Recall}
\begin{itemize}
  \item OS = application controller + hardware interface + resource manager.
  \item OS hides hardware details behind services and interfaces.
  \item Important services: development, execution, I/O, files, access control, errors, accounting.
  \item Kernel holds common privileged OS functions in memory.
  \item System calls are the legal bridge from user mode to kernel mode.
  \item Evolution needs modular design, clear interfaces, and maintainability.
\end{itemize}

\cheatsection{3 Process Description and Control}
\cheatsubsection{Process, Image, PCB}
\term[进程/执行中的程序]{Process}: a program in execution; an entity that can be assigned to and executed by a processor.\\
\term[进程映像]{Process image}: program code, data, stack, and process attributes.\\
\term[进程控制块]{PCB}: OS record holding enough information to stop and later resume a process as if it had not been interrupted.
\begin{tabularx}{\linewidth}{@{}TY@{}}
\toprule
\term[标识信息]{ID} & PID, parent PID, user ID.\\
\term[处理器状态]{CPU state} & registers, PC, PSW, stack pointer.\\
\term[控制信息]{Control} & state, priority, scheduling links, memory pointers, I/O status, accounting.\\
\bottomrule
\end{tabularx}
\term[指令轨迹]{Instruction trace}: sequence of instructions actually executed by a process.\\
\term[OS管理表]{OS tables}: memory, I/O, file, and process tables.

\cheatsubsection{Five-State Model}
\begin{tabularx}{\linewidth}{@{}TY@{}}
\toprule
\term[新建]{New} & created but not yet admitted.\\
\term[就绪/只差CPU]{Ready} & can run once assigned a CPU; lacks only CPU.\\
\term[运行]{Running} & currently executing on the CPU.\\
\term[阻塞/等事件]{Blocked} & waits for an event such as I/O completion.\\
\term[退出]{Exit} & terminated or released.\\
\bottomrule
\end{tabularx}
\term[接纳]{Admit}: New \(\to\) Ready. \term[派遣]{Dispatch}: Ready \(\to\) Running. \term[超时/抢占]{Timeout/preempt}: Running \(\to\) Ready. \term[等待事件]{Event wait}: Running \(\to\) Blocked. \term[事件发生]{Event occurs}: Blocked \(\to\) Ready. \term[释放]{Release}: Running \(\to\) Exit.\\
\note{Ready = waiting for CPU; Blocked = waiting for an event. An event-completed blocked process returns to Ready, not directly to Running.}

\cheatsubsection{Suspension and Swapping}
\term[交换/内存磁盘搬移]{Swapping}: moving a process between main memory and disk to free memory and improve CPU use.\\
\term[挂起/不可立即执行]{Suspended}: not immediately available for execution; may or may not be waiting for an event.
\begin{tabularx}{\linewidth}{@{}TY@{}}
\toprule
\term[就绪挂起]{Ready/Suspend} & swapped out; event condition satisfied; can run after swap-in.\\
\term[阻塞挂起]{Blocked/Suspend} & swapped out and still waiting for an event.\\
\term[挂起原因]{Reasons} & swapping, OS reason, interactive user request, timing, parent request.\\
\term[挂起特征]{Traits} & not executable now; may wait on event; agent-suspended; needs explicit activation.\\
\bottomrule
\end{tabularx}
\term[队列组织]{Queues}: use one Ready queue plus per-event Blocked queues; suspended states use Ready/Suspend and per-event Blocked/Suspend queues.\par
\cheatsubsection{Seven-State Diagram}
\noindent\includegraphics[width=\linewidth]{chapter3-seven-state.png}

\cheatsubsection{Control, Modes, Switching}
\term[创建进程步骤]{Creation steps}: assign PID; allocate process space; initialize PCB; link into ready or suspend queues; create/expand accounting and other tables.\\
\term[用户态]{User mode}: low privilege; cannot execute privileged instructions. \term[内核/系统态]{Kernel/System mode}: high privilege for protected OS operations.\\
\term[外部异步]{Interrupt}: asynchronous external event. \term[当前指令异常]{Trap}: exception caused by the current instruction. \term[访管调用]{Supervisor call}: explicit request for an OS service.\\
\term[模式切换]{Mode switch}: changes privilege mode, possibly same process. \term[进程切换]{Process switch}: CPU changes from one process to another.\\
\term[上下文切换工作]{Process switch work}: save old PC, PSW, registers into PCB; update state and queues; choose next process; load next PCB/context.

\cheatsection{Chapter 3 Review Q\&A}
\begin{stemlist}
\stemrow[指令轨迹]{R3.1 Trace}{Stem: instruction trace. A: executed instruction sequence of a process.}
\stemrow[进程创建来源]{R3.2 Creation}{Stem: process creation events. A: new batch job, interactive logon, OS-created service, spawned child.}
\stemrow[五状态]{R3.3 States}{Stem: Fig. 3.6 states. A: New, Ready, Running, Blocked, Exit.}
\stemrow[抢占定义]{R3.4 Preempt}{Stem: preempt a process. A: reclaim CPU before the process voluntarily yields or finishes.}
\stemrow[交换目的]{R3.5 Swapping}{Stem: swapping and purpose. A: move process memory to/from disk to free memory and improve CPU use.}
\stemrow[内存内外阻塞]{R3.6 Two blocked}{Stem: why Fig. 3.9b has two blocked states. A: resident Blocked vs swapped-out Blocked/Suspend.}
\stemrow[挂起特征]{R3.7 Suspended}{Stem: four suspended-process traits. A: not executable now; may wait; agent-suspended; explicitly activated.}
\stemrow[OS四类表]{R3.8 OS tables}{Stem: OS management tables. A: memory, I/O, file, process.}
\stemrow[PCB三类信息]{R3.9 PCB info}{Stem: three PCB categories. A: identification, processor state, process control.}
\stemrow[模式保护]{R3.10 Modes}{Stem: why user/kernel modes. A: protect OS, PCBs, memory, I/O, privileged instructions.}
\stemrow[创建步骤]{R3.11 Create}{Stem: OS steps to create process. A: PID, space, PCB, links, other tables.}
\stemrow[中断vs陷阱]{R3.12 Int/trap}{Stem: interrupt vs trap. A: external asynchronous event vs instruction-caused exception.}
\stemrow[中断例子]{R3.13 Examples}{Stem: three interrupts. A: clock, I/O, memory fault.}
\stemrow[模式vs进程切换]{R3.14 Switches}{Stem: mode switch vs process switch. A: privilege change vs changing the running process.}
\end{stemlist}

\cheatsection{Chapter 3 Problem Q\&A}
\begin{stemlist}
\stemrow[状态转换表]{P3.1 Transitions}{Stem: R/RUN/B/NR transition table. A: 1 dispatch; 2 preempt; 3 event wait; 4 event occurs; 5 ready swap-out; 6 blocked swap-out.}
\stemrow[按时间判状态]{P3.2 Timeline}{Stem: process states at \(t=22,37,47\). A: \(t=22\): P1 B(d3 read), P3 B(d2 read), P7 B(d3 write), P5 likely running. \(t=37\): P1/P3 ready, P5 Blocked/Suspend, P7 blocked, P8 running. \(t=47\): P5 ready, P7 blocked, P8 exit; one ready process may run.}
\stemrow[七状态可达性]{P3.3 Seven states}{Stem: possible/impossible Fig. 3.9b transitions. A: Possible direct moves: N\(\to\)R/RS/E; R\(\to\)Run/RS/E; Run\(\to\)R/B/RS/E; B\(\to\)R/BS/E; RS\(\to\)R/E; BS\(\to\)B/RS/E; all other direct moves are impossible.}
\stemrow[就绪/阻塞队列]{P3.4 Queues}{Stem: seven-state queueing diagram. A: Ready queue, processor, per-event Blocked queues, Ready/Suspend queue, and per-event Blocked/Suspend queues.}
\stemrow[换入策略]{P3.5 Policy}{Stem: choose Ready vs high-priority Ready/Suspend. A: swap in Ready/Suspend only when priority or aging benefit exceeds swap cost; otherwise dispatch resident Ready.}
\stemrow[VAX等待状态]{P3.6 VAX states}{Stem: justify VAX/VMS wait states. A: many waits give precise event/resource queues; paging/global resource waits need no resident/outswapped pair; transitions follow dispatch, wait, event, swap.}
\stemrow[多访问模式]{P3.7 Access modes}{Stem: four access modes vs two. A: Kernel, Executive, Supervisor, User give finer least-privilege protection, but add complexity and mode-switch overhead.}
\stemrow[环保护限制]{P3.8 Rings}{Stem: ring protection need-to-know problem. A: rings are nested; inner domains can access outer objects, so selective need-to-know access is hard to enforce.}
\stemrow[多事件等待]{P3.9 Multi-wait}{Stem: process waiting on multiple events. A: yes, e.g., network, keyboard, timeout; use wait descriptors or multiple queue links and remove stale links after wakeup.}
\stemrow[固定保存位置]{P3.10 Fixed saves}{Stem: fixed interrupt register-save locations. A: practical only for simple nonnested systems; nested interrupts, multiprogramming, and switches need PCB/kernel-stack context.}
\stemrow[非抢占内核风险]{P3.11 Real-time}{Stem: nonpreemptive kernel and real-time. A: it can cause unbounded latency, priority inversion, and missed deadlines; real-time systems need bounded response.}
\stemrow[fork输出顺序]{P3.12 Fork}{Stem: possible output of successful \code{fork()}. A: two lines print: \(0\) in the child and child PID in the parent; order is scheduler-dependent.}
\end{stemlist}

\cheatsection{Chapter 3 Fast Recall}
\begin{itemize}
  \item Process = program in execution; PCB makes pause/resume possible.
  \item Ready lacks CPU; Blocked lacks an event; Suspended is out of main memory.
  \item Preemption takes CPU from a process before it yields.
  \item Interrupts are external; traps are instruction-caused; supervisor calls request OS services.
  \item A mode switch can stay in the same process; a process switch always changes processes.
  \item Context switching saves old PCB state and restores the selected process state.
\end{itemize}



\cheatsection{4 Threads}
\cheatsubsection{Process vs Thread}
\term[资源拥有单位]{Process}: unit of resource ownership and protection; owns address space, process image, open files, I/O resources, privileges, and IPC links.\\
\term[调度执行单位]{Thread}: unit of scheduling/execution; has state, PC, registers, stack, priority, context, and thread-local data. Threads in one process share memory and resources.\\
\term[多线程]{Multithreading}: multiple concurrent execution paths inside one process. \term[进程资源]{PCB} stores process resources; \term[线程上下文]{TCB} stores thread execution context.\\
\term[线程优点]{Advantages}: faster create/terminate/switch, cheaper communication, better responsiveness, overlapped I/O and computation, true parallelism on multiprocessors.

\cheatsubsection{States, Operations, Synchronization}
\term[线程状态]{Thread states}: Running, Ready, Blocked. Suspend states are usually process-level because all threads share the process address space.\\
\term[线程操作]{Operations}: spawn creates a ready thread; block waits for an event; unblock moves to Ready; finish releases context and stack.\\
\term[同步/防共享冲突]{Synchronization}: required because shared address space/resources can cause data races, lost updates, and corrupted data structures.

\cheatsubsection{ULT, KLT, Combined}
\term[用户级线程]{ULT}: managed by a user-level thread library; kernel sees only the process. Pros: no kernel-mode switch, application scheduling, portable. Cons: blocking syscall may block all process threads; pure ULT cannot exploit multiprocessing.\\
\term[内核级线程]{KLT}: managed and scheduled by the kernel. Pros: one blocked thread need not block the process; threads can run on multiple CPUs; kernel code can be multithreaded. Cons: higher switch overhead.\\
\term[封套/包裹阻塞调用]{Jacketing}: wrap a blocking system call so the thread library can test/yield and avoid blocking the whole process.\\
\term[混合模型]{Combined}: map many ULTs onto some KLTs/LWPs to mix fast user-level management with kernel parallelism.

\cheatsection{Chapter 4 Review Q\&A}
\begin{stemlist}
\stemrow[PCB和TCB分工]{R4.1 PCB/TCB}{Stem: which PCB elements move to TCB. A: TCB gets thread state, PC, registers, stack pointer, priority, scheduling info, private data; PCB keeps PID, address space, memory, files/I/O, privileges, IPC.}
\stemrow[线程切换便宜]{R4.2 Switch cost}{Stem: why thread mode switch may be cheaper. A: same-process threads keep the same address space and resources; only thread context is saved/restored.}
\stemrow[进程两属性]{R4.3 Process traits}{Stem: two independent process characteristics. A: resource ownership and scheduling/execution.}
\stemrow[线程用途]{R4.4 Uses}{Stem: four thread uses. A: foreground/background work, asynchronous processing, faster execution by overlap/parallelism, modular program structure.}
\stemrow[共享/私有资源]{R4.5 Shared resources}{Stem: resources shared by process threads. A: address space, code, global data, heap, files, I/O resources, privileges; each thread has own PC, registers, stack, state.}
\stemrow[ULT优点]{R4.6 ULT pros}{Stem: three ULT advantages over KLTs. A: no kernel switch for thread ops, application-specific scheduling, portability.}
\stemrow[ULT缺点]{R4.7 ULT cons}{Stem: two ULT disadvantages. A: blocking syscall may block the whole process; pure ULT cannot use multiple CPUs in parallel.}
\stemrow[封套定义]{R4.8 Jacketing}{Stem: define jacketing. A: wrapper turns a blocking syscall into nonblocking behavior for user-level thread scheduling.}
\end{stemlist}

\cheatsection{Chapter 4 Problem Q\&A}
\begin{stemlist}
\stemrow[同进程切换]{P4.1 Same-process switch}{Stem: same-process vs different-process thread switch. A: yes; same-process switch avoids address-space/resource switch.}
\stemrow[ULT阻塞全进程]{P4.2 ULT blocking}{Stem: why one blocking ULT blocks all. A: kernel is unaware of ULTs and blocks the whole process on a blocking syscall.}
\stemrow[会话与资源归属]{P4.3 OS/2 sessions}{Stem: remove sessions, bind UI to processes. A: lose grouping of several processes under one UI; assign memory/files/resources to processes, not threads.}
\stemrow[I/O重叠]{P4.4 1:1 on uniprocessor}{Stem: why faster than single thread. A: when one thread waits for I/O, another ready thread can run; CPU use and responsiveness improve.}
\stemrow[进程结束线程全灭]{P4.5 Process exit}{Stem: process exits while threads run. A: no thread continues; threads depend on the process address space/resources.}
\stemrow[地址空间/任务分离]{P4.6 OS/390 ASCB/TCB}{Stem: split global ASCB and local TCB. A: separates address-space management from task execution and reduces global kernel data.}
\stemrow[共享计数竞态]{P4.7 count\_positives}{Stem: function and two-thread risk. A: counts positive list values into a global counter; unsynchronized shared counter creates data-race risk.}
\stemrow[寄存器缓存丢更新]{P4.8 Optimized code}{Stem: compiler caches global counter. A: local-register copy can overwrite another thread update, causing lost update.}
\stemrow[自增非原子]{P4.9 Pthreads}{Stem: two threads increment \code{myglobal}; output 21. A: unexpected race; read-add-write is not atomic, so increments are lost.}
\stemrow[yield同优先级]{P4.10 Solaris yield}{Stem: yield to same priority only. A: a higher-priority runnable thread would normally be selected by the scheduler; yield donates only among equal-priority peers.}
\stemrow[ULT到LWP映射]{P4.11 Solaris ULT/LWP}{Stem: many ULTs per LWP and state diagrams. A: fewer LWPs reduce kernel overhead; ULT state and LWP state differ because user library and kernel schedule different entities.}
\stemrow[不可中断等待]{P4.12 Linux uninterruptible}{Stem: rationale for Uninterruptible state. A: protect kernel/device consistency while waiting for critical I/O/resource events; wake only when event completes.}
\end{stemlist}

\cheatsection{Chapter 4 Fast Recall}
\begin{itemize}
  \item Process owns resources; thread executes code.
  \item Threads share process memory/resources but keep private PC, registers, stack, and state.
  \item ULT is fast and portable but weak on blocking and multiprocessing.
  \item KLT costs more but supports kernel scheduling, blocking isolation, and CPU parallelism.
  \item Race condition = unsynchronized shared read/write; lost update is the classic symptom.
\end{itemize}

\cheatsection{5 Concurrency and Semaphores}
\cheatsubsection{Concurrency Basics}
\term[并发/同一时间段推进]{Concurrency}: multiple processes or threads make progress in the same time interval. On one CPU execution is interleaved; on many CPUs it may overlap in parallel.\\
\term[并发风险]{Main hazards}: shared global resources, hard resource allocation, and nondeterministic bugs(program hard to reproduce).\\
\term[竞态/结果看时序]{Race condition}: final result depends on timing of unsynchronized shared read/write.\\
\term[临界区/访问共享资源]{Critical section}: code that accesses a shared resource. \term[互斥/一次一个]{Mutual exclusion}: at most one process may execute the corresponding critical section for the same resource.

\cheatsubsection{Interaction and OS Concerns}
\begin{tabularx}{\linewidth}{@{}TY@{}}
\toprule
\term[互不知道/竞争]{Unaware} & competition; issues: mutual exclusion, deadlock, starvation.\\
\term[间接知道/共享数据]{Indirectly aware} & cooperation by sharing; add data coherence risk.\\
\term[直接知道/通信]{Directly aware} & cooperation by communication; no shared-data mutex, but deadlock/starvation still possible.\\
\bottomrule
\end{tabularx}
\term[OS并发职责]{OS duties}: track processes, allocate/deallocate resources, protect data/resources, and ensure speed independence.

\cheatsubsection{Mutual Exclusion Requirements}
\begin{itemize}
  \item Only one process in a critical section for a resource.
  \item A halted noncritical process must not block others.
  \item No deadlock or starvation.
  \item If the critical section is free, entry is not delayed.
  \item Assume no process speed or processor count.
  \item A process stays in the critical section only finite time.
\end{itemize}

\cheatsubsection{Semaphores}
\term[信号量/原子P V]{Semaphore}: integer synchronization variable accessed only by atomic operations: initialize, \code{semWait}/P/down, \code{semSignal}/V/up.\\
\term[计数含义]{Count meaning}: \(s>0\) means available units; \(s=0\) means none free; \(s<0\) means \(|s|\) blocked waiters.\\
\term[二元信号量]{Binary semaphore}: value 0/1, often mutex-like. \term[计数信号量]{General/counting semaphore}: integer count for multiple resources or ordering.\\
\term[互斥锁vs二元信号量]{Mutex vs binary semaphore}: a mutex is unlocked by its locker; a binary semaphore may be waited by one process and signaled by another.\\
\term[强信号量/FIFO]{Strong semaphore}: FIFO wait queue. \term[弱信号量/可能饥饿]{Weak semaphore}: no specified release order; starvation is possible.\\
\term[实现要点/原子性]{Implementation}: \code{semWait} and \code{semSignal} must be atomic; use hardware support, compare-and-swap, short interrupt disable on uniprocessors, or software algorithms.

\cheatsubsection{Producer/Consumer Patterns}
\term[无限缓冲/只数产品]{Infinite buffer}: \code{n=0} counts items; \code{s=1} protects the buffer. Producer: wait \(s\), append, signal \(s\), signal \(n\). Consumer: wait \(n\), wait \(s\), take, signal \(s\).\\
\term[有界缓冲/空位+产品]{Bounded buffer}: \code{e=sizeofbuffer}, \code{n=0}, \code{s=1}. Producer: wait \(e\), wait \(s\), append, signal \(s\), signal \(n\). Consumer: wait \(n\), wait \(s\), take, signal \(s\), signal \(e\).\\
\note{Never wait on the resource count while holding the buffer mutex: producer wait \(s\) then \(e\), or consumer wait \(s\) then \(n\), can deadlock.}

\cheatsection{Chapter 5 Review Q\&A}
\begin{stemlist}
\stemrow[并发四问题]{R5.1 Design issues}{Stem: four concurrency design issues. A: process communication, resource sharing/competition, synchronization, processor-time allocation.}
\stemrow[并发出现位置]{R5.2 Contexts}{Stem: three contexts where concurrency arises. A: multiple applications, structured applications, OS structure.}
\stemrow[基本要求互斥]{R5.3 Basic requirement}{Stem: requirement for concurrent execution. A: enforce mutual exclusion for critical resources, with correct synchronization.}
\stemrow[感知三层次]{R5.4 Awareness}{Stem: three degrees of process awareness. A: unaware/competition; indirectly aware/cooperation by sharing; directly aware/cooperation by communication.}
\stemrow[竞争vs协作]{R5.5 Compete vs cooperate}{Stem: distinction. A: competing processes only contend for resources; cooperating processes work together by shared data or messages.}
\stemrow[互斥死锁饥饿]{R5.6 Control problems}{Stem: three competing-process problems. A: mutual exclusion, deadlock, starvation.}
\stemrow[互斥六条件]{R5.7 Mutex requirements}{Stem: list requirements. A: one-at-a-time, noncritical halt harmless, no deadlock/starvation, no needless delay, no speed/CPU assumptions, finite CS time.}
\stemrow[信号量操作]{R5.8 Semaphore ops}{Stem: operations. A: initialize, \code{semWait}/P, \code{semSignal}/V; no direct arbitrary access.}
\stemrow[二元vs计数]{R5.9 Binary/general}{Stem: difference. A: binary is 0/1; general/counting is integer-valued for multiple units or synchronization.}
\stemrow[强弱信号量]{R5.10 Strong/weak}{Stem: difference. A: strong uses FIFO waiter release; weak gives no release-order guarantee.}
\stemrow[管程定义]{R5.11 Monitor}{Stem: monitor. A: high-level synchronization module encapsulating shared data/procedures; only one process active inside, with condition variables for waiting.}
\stemrow[消息阻塞性]{R5.12 Messages}{Stem: blocking vs nonblocking. A: blocking send/receive waits for completion or message; nonblocking returns immediately after queueing or status check.}
\stemrow[读者写者]{R5.13 Readers/writers}{Stem: conditions. A: many readers may read together; writers need exclusive access; policies must avoid reader or writer starvation.}
\end{stemlist}

\cheatsection{Chapter 5 Problem Q\&A}
\begin{stemlist}
\stemrow[单CPU vs 多CPU]{P5.1 Multi vs multiprocessing}{Stem: two concurrency differences. A: multiprogramming interleaves one CPU and can protect short regions by disabling interrupts; multiprocessing has true overlap, so one-CPU interrupt disable is insufficient.}
\stemrow[竞态交错]{P5.3 Race traces}{Stem: print \code{x is 10} or \code{x is 8}. A: interleave the \code{if} test and \code{printf} to print 10; split load/decrement/store and load/increment/store to overwrite with stale 8.}
\stemrow[丢更新上下界]{P5.4 Tally bounds}{Stem: two \code{total()} processes, \(n=50\). A: final \code{tally} is 2 to 100; with \(k\) processes, upper \(50k\), lower can still be 2 for \(k\ge2\).}
\stemrow[忙等适用场景]{P5.5 Busy wait}{Stem: always less efficient than blocking? A: no; short waits may be cheaper than block/context-switch/restart overhead; long waits should block.}
\stemrow[信号量等价定义]{P5.12 Sem definitions}{Stem: negative-count vs nonnegative-count semaphore definitions. A: equivalent to programs because waiters are represented either by negative count or by the queue; users cannot inspect internals.}
\stemrow[二元实现计数]{P5.16 General via binary}{Stem: flaw and fix. A: binary \code{delay} can lose wakeups; use pass-the-baton: if \code{semSignal} wakes a waiter, it signals \code{delay} instead of releasing \code{mutex}; awakened waiter releases \code{mutex}.}
\stemrow[拿锁等待导致死锁]{P5.19 Consumer deadlock}{Stem: move empty-test inside mutex. A: consumer can hold \(s\) while waiting on \code{delay}; producer then waits for \(s\), so both block.}
\stemrow[减少信号量开销]{P5.20 Faster infinite buffer}{Stem: reduce semaphore overhead. A: let \(n=-1\) mean empty and consumer waiting; producer signals \code{delay} only when append changes \(n\) from -1 to 0.}
\stemrow[wait顺序会死锁]{P5.21 Swap order}{Stem: bounded-buffer statement swaps. A: swapping wait \(e/s\) or wait \(n/s\) can deadlock; swapping signal \(s/n\) or signal \(s/e\) preserves correctness but may add waiting.}
\stemrow[有界缓冲追踪]{P5.22 Trace table}{Stem: fill bounded-buffer chart. A: blank labels: Pa p2, Pb p2, Pb p3, Cb c2, Pc p2, Cb c3, Pa p3, Pc p3, Pa p2, Pa p3; use FIFO queues; final state full \(=3\), empty \(=0\), buffer \(=XXX\).}
\end{stemlist}

\cheatsection{Chapter 5 Fast Recall}
\begin{itemize}
  \item Race condition = timing-dependent result on shared data.
  \item Protect the code that touches shared state, not just the variable name.
  \item P/down may block; V/up may wake zero or one waiter.
  \item Semaphore waits before mutex waits in producer/consumer resource tests.
  \item Strong semaphore fairness comes from FIFO; weak semaphore may starve.
\end{itemize}

\cheatsection{6 Deadlock and Starvation}
\cheatsubsection{Resources and Deadlock Conditions}
\term[可重用资源]{Reusable resource}: not depleted by use; released and used again, e.g., CPU, memory, files, devices, semaphores.\\
\term[可消耗资源]{Consumable resource}: produced and consumed; after use it is gone, e.g., messages, signals, interrupts, buffer information.\\
\term[死锁/互等永久阻塞]{Deadlock}: processes are permanently blocked because each waits for an event/resource held by another process. \term[饥饿/一直不被选]{Starvation}: a ready/requesting process waits indefinitely because others keep being favored.
\begin{tabularx}{\linewidth}{@{}TY@{}}
\toprule
\term[互斥条件]{Mutual exclusion} & resource can be held by only one process.\\
\term[占有并等待]{Hold and wait} & process holds resources while requesting more.\\
\term[不可抢占]{No preemption} & resource cannot be forcibly taken away.\\
\term[循环等待]{Circular wait} & closed chain: each process waits for a resource held by the next.\\
\bottomrule
\end{tabularx}
\note{First three conditions make deadlock possible; adding circular wait creates deadlock.}

\cheatsubsection{Handling Strategies}
\term[预防/破坏条件]{Prevention}: make at least one deadlock condition impossible. Examples: request all resources at once, allow preemption where safe, or impose a global resource order.\\
\term[避免/只给安全状态]{Avoidance}: grant a request only if the resulting state is safe. Needs each process's maximum claim.\\
\term[检测/事后发现恢复]{Detection}: grant resources more freely, periodically test the current state, then recover by aborting, rollback, or preemption.

\cheatsubsection{Banker's Algorithm}
\[
\mathrm{Need}=\mathrm{Claim}-\mathrm{Allocation}
\]
For request \(R_i\): require \(R_i\le Need_i\) and \(R_i\le Available\); pretend to allocate; grant only if a safe sequence exists.\\
\term[安全状态/有完成序列]{Safe state}: some process completion order can satisfy all remaining needs. \term[不安全/未必已死锁]{Unsafe state}: no guaranteed safe sequence; deadlock may occur but has not necessarily occurred.

\cheatsubsection{Avoidance vs Detection}
\begin{tabularx}{\linewidth}{@{}TY@{}}
\toprule
\term[避免算法]{Avoidance} & before allocation; uses Claim/Max, Allocation, Available; tests \(Need_i\le Available\); failure means unsafe.\\
\term[检测算法]{Detection} & after allocation; uses Request \(Q\), Allocation, Available; tests \(Q_i\le W\); unmarked at end means deadlocked.\\
\bottomrule
\end{tabularx}
\term[检测步骤]{Detection sketch}: set \(W=Available\); mark processes with zero allocation; repeatedly find unmarked \(Q_i\le W\), mark it, set \(W=W+A_i\). Remaining unmarked processes are deadlocked.

\cheatsection{Chapter 6 Review Q\&A}
\begin{stemlist}
\stemrow[资源类型例子]{R6.1 Resource examples}{Stem: reusable vs consumable resources. A: reusable: CPU, memory, files, devices, semaphores; consumable: messages, signals, interrupts, buffer data.}
\stemrow[死锁可能三条件]{R6.2 Possible deadlock}{Stem: three conditions for deadlock possibility. A: mutual exclusion, hold and wait, no preemption.}
\stemrow[死锁四条件]{R6.3 Actual deadlock}{Stem: four deadlock conditions. A: mutual exclusion, hold and wait, no preemption, circular wait.}
\stemrow[破坏占有等待]{R6.4 Prevent hold-wait}{Stem: hold-and-wait prevention. A: require a process to request all needed resources at once, or release held resources before requesting more.}
\stemrow[破坏不可抢占]{R6.5 Prevent no-preempt}{Stem: two no-preemption attacks. A: if a request fails, release held resources and retry; or preempt resources from another holder when state can be restored.}
\stemrow[资源全序]{R6.6 Prevent circular wait}{Stem: circular-wait prevention. A: impose a total order on resource types and request only in increasing order.}
\stemrow[预防/避免/检测]{R6.7 Three strategies}{Stem: prevention vs avoidance vs detection. A: prevention breaks a condition; avoidance refuses unsafe grants; detection finds existing deadlocks and recovers.}
\end{stemlist}

\cheatsection{Chapter 6 Problem Q\&A}
\begin{stemlist}
\stemrow[交通四条件]{P6.1 Traffic deadlock}{Stem: show four conditions in Fig. 6.1a. A: each road quadrant is exclusive; cars hold one quadrant while waiting; cars cannot be preempted; cars wait in a cycle.}
\stemrow[三策略套交通]{P6.2 Apply strategies}{Stem: prevention/avoidance/detection for Fig. 6.1. A: traffic lights/order break circular wait; controller allows only safe entry; detection notices blocked cycle and backs/removes one car.}
\stemrow[路径避死锁]{P6.3 Six paths}{Stem: narrate Fig. 6.3 paths. A: Q-first, P-then-Q, mixed Q-holds-B/P-releases-A, mixed after A release, P-holds-B/Q-waits, or P-first; all avoid a fatal region.}
\stemrow[无循环等待]{P6.4 No deadlock}{Stem: justify Fig. 6.3. A: P releases A before requesting B; Q may hold B and wait for A, but P does not hold A while waiting for B, so no circular wait.}
\stemrow[银行家安全序列]{P6.5 Banker state}{Stem: verify, need, safe sequence, P5 request. A: allocated total 9346, so Available 6354. Need rows: P0 7534, P1 2122, P2 3442, P3 2331, P4 4121, P5 3433. Safe order: P1,P2,P0,P3,P4,P5. Granting P5 3233 leaves 3121 and no process can finish; deny.}
\par\vspace{1pt}\noindent
\begin{tabular}{@{}c@{\hspace{1pt}}c@{}}
\includegraphics[width=0.49\linewidth]{chapter6-rag-requested.png} &
\includegraphics[width=0.49\linewidth]{chapter6-rag-held.png}\\[-1pt]
\includegraphics[width=0.49\linewidth]{chapter6-rag-circular-wait.png} &
\includegraphics[width=0.49\linewidth]{chapter6-rag-no-deadlock.png}
\end{tabular}
\par\vspace{1pt}
\stemrow[资源分配图+排序]{P6.6 RAG/order}{Stem: code with P0 A,B,C; P1 D,E,B; P2 C,F,D. A: deadlock cycle P0 waits C held by P2; P2 waits D held by P1; P1 waits B held by P0. Prevent by global order A<B<C<D<E<F: P0 A,B,C; P1 B,D,E; P2 C,D,F.}
\stemrow[内存版银行家]{P6.11 Memory banker}{Stem: total 150, P1 70/45, P2 60/40, P3 60/15; add P4. A: initial Available \(=50\). P4 hold 25 gives Available 25 and safe order P1,P2,P3,P4; grant. P4 hold 35 gives Available 15 and no need fits; deny.}
\stemrow[银行家局限]{P6.12 Banker usefulness}{Stem: evaluate banker's algorithm in an OS. A: good theory and predictable-resource systems; often impractical because max claims must be known, resources/processes vary, each request costs a safety test, and concurrency may be reduced.}
\stemrow[最小Available]{P6.15 Min available}{Stem: one resource, \(C=(3,2,9,7)\), \(A=(1,1,3,2)\). A: Need \(=(2,1,6,5)\); minimum Available is 3, e.g., P2,P1,P4,P3.}
\stemrow[方法效率排序]{P6.16 Rankings}{Stem: compare six deadlock methods by concurrency and efficiency. A: concurrency high to low: detect+kill, detect+rollback, banker, ordering, restart-on-wait, reserve-all. Rare-deadlock efficiency: ordering, reserve-all, detect+kill, restart-on-wait, banker, detect+rollback; if deadlocks are frequent, detection recovery becomes less attractive.}
\end{stemlist}

\cheatsection{Chapter 6 Fast Recall}
\begin{itemize}
  \item Prevention breaks conditions; avoidance checks future safety; detection checks current deadlock.
  \item Unsafe does not mean deadlocked; it means no guaranteed safe completion order.
  \item Banker's core loop: find \(Need_i\le Available\), finish \(P_i\), then add back \(Allocation_i\).
  \item Detection core loop uses current requests \(Q_i\), not maximum claims.
  \item Global resource ordering prevents circular wait but can reduce flexibility.
\end{itemize}

\cheatsection{7 Memory Management}
\cheatsubsection{Requirements and Addresses}
\term[内存管理]{Memory management}: allocates main memory among active processes and supports efficient multiprogramming.\\
\term[五个需求]{Requirements}: relocation, protection, sharing, logical organization, and physical organization.\\
\term[逻辑地址/程序看到]{Logical address}: program-visible address. \term[相对地址/偏移]{Relative address}: logical offset from a known point. \term[物理地址/真实主存]{Physical address}: real main-memory address.\\
\term[动态重定位]{Dynamic relocation}: hardware maps relative to physical address, commonly \(physical=base+relative\); bounds/limit checks enforce protection.

\cheatsubsection{Partitioning}
\term[固定分区/内碎片]{Fixed partitioning}: system-generation split; simple, but internal fragmentation and fixed process count.\\
\term[动态分区/外碎片]{Dynamic partitioning}: exact-size allocation; no internal fragmentation, but external fragmentation and possible compaction.\\
\term[首次适配]{First-fit}: first large enough hole; fast and usually good. \term[下次适配]{Next-fit}: resumes from last placement; often uses later holes.\\
\term[最佳适配]{Best-fit}: smallest large enough hole; often leaves many tiny holes.\\
\term[内部碎片]{Internal frag.}: waste inside allocation. \term[外部碎片]{External frag.}: separated holes.

\cheatsubsection{Buddy System and Paging}
\term[伙伴系统/2的幂]{Buddy system}: block sizes are powers of two; split on allocation and coalesce free buddies on release. Buddy address: \(buddy_k(x)=x\oplus2^k\).\\
\term[分页/页框等大]{Paging}: process logical space is split into fixed-size pages; physical memory is split into same-size frames. Pages need not be contiguous.\\
\term[页表/页到框]{Page table}: maps page number to frame number. Logical address \(=(page,offset)\); physical address \(=(frame,offset)\).
\[
\begin{aligned}
page&=\lfloor LA/page\_size\rfloor\\
offset&=LA\bmod page\_size\\
frame&=page\_table[page]\\
PA&=frame\times page\_size+offset
\end{aligned}
\]
\note{Paging has no external fragmentation; only the last page may have internal fragmentation.}

\cheatsection{Chapter 7 Review Q\&A}
\begin{stemlist}
\stemrow[五大需求]{R7.1 Requirements}{Stem: memory-management requirements. A: relocation, protection, sharing, logical organization, physical organization.}
\stemrow[重定位好处]{R7.2 Relocation}{Stem: why relocation is desirable. A: processes can be loaded, swapped, or compacted into different physical locations without changing program addresses.}
\stemrow[运行时保护]{R7.3 Compile protection}{Stem: why protection cannot be enforced at compile time. A: final physical locations, sharing, and relocation are run-time facts, so run-time checks are needed.}
\stemrow[共享用途]{R7.4 Sharing}{Stem: why share memory regions. A: shared libraries/code, shared data, IPC, common files/buffers, and memory savings.}
\stemrow[不等分区优点]{R7.5 Unequal fixed partitions}{Stem: advantages. A: better match to process sizes, less internal waste, and support for both small and larger processes.}
\stemrow[内碎片vs外碎片]{R7.6 Fragmentation}{Stem: internal vs external. A: internal is unused space inside an allocation; external is separated free holes outside allocations.}
\stemrow[三类地址]{R7.7 Address types}{Stem: logical/relative/physical. A: logical is program-visible; relative is an offset; physical is the real main-memory location.}
\stemrow[页vs页框]{R7.8 Page/frame}{Stem: difference. A: a page is a fixed-size logical-process block; a frame is the same-size physical-memory block.}
\stemrow[分页vs分段]{R7.9 Page/segment}{Stem: difference. A: pages are fixed-size and system-managed; segments are variable-size logical units such as code, data, or stack.}
\end{stemlist}

\cheatsection{Chapter 7 Problem Q\&A}
\begin{stemlist}
\stemrow[目标对照需求]{P7.1 Objectives}{Stem: Section 2.3 objectives vs Section 7.1 requirements. A: isolation/protection, allocation/relocation, modularity/logical organization, access control/sharing, and storage hierarchy/physical organization cover the same concerns.}
\stemrow[分区数取log]{P7.2 Partition pointer}{Stem: \(2^{16}\)-byte partitions in \(2^{24}\)-byte memory. A: \(2^{24}/2^{16}=2^8\) partitions, so pointer needs 8 bits.}
\stemrow[50百分比规则]{P7.3 50 percent rule}{Stem: holes vs segments in dynamic partitioning. A: in equilibrium, random deletion/allocation gives about one hole per two allocated segments, so holes \(\approx n/2\).}
\stemrow[空闲表搜索]{P7.4 Search length}{Stem: average free-list search. A: best-fit checks all holes unless size-sorted; first-fit and next-fit scan roughly half the list on average.}
\stemrow[最坏适配]{P7.5 Worst-fit}{Stem: largest hole placement. A: may avoid tiny leftovers but destroys large holes; search is full list unless kept largest-first.}
\stemrow[分区算法图]{P7.6 Diagram}{Stem: X is last 2 MB placement; then one swap-out. A: under low-address splitting, swapped-out max \(=1\) MB; X split a 7 MB hole; next 3 MB: best-fit exact 3 MB hole, first-fit first 4 MB hole, next-fit 5 MB hole after X, worst-fit 8 MB hole.}
\stemrow[伙伴分裂合并]{P7.7 Buddy sequence}{Stem: 1 MiB buddy requests 70,35,80; returns A; request 60; returns B,D,C. A: A=128 KiB, B=64 KiB, C=128 KiB, D=64 KiB; after Return B: free A128, free B64, used D64, used C128, free 128, free 512; final all coalesces to 1 MiB free.}
\stemrow[伙伴地址翻位]{P7.8 Buddy address}{Stem: block address \code{011011110000}. A: size 4 buddy \code{011011110100}; size 16 buddy \code{011011100000}.}
\stemrow[异或公式]{P7.9 Buddy formula}{Stem: general buddy expression. A: \(buddy_k(x)=x\oplus2^k\); equivalently flip bit \(k\).}
\stemrow[斐波那契伙伴]{P7.10 Fibonacci buddy}{Stem: use Fibonacci sizes. A: yes; split \(F_{n+2}\) into \(F_{n+1}\) and \(F_n\). Advantage: less internal fragmentation; cost: harder buddy management.}
\stemrow[PC用相对地址]{P7.11 PC address}{Stem: relative vs absolute PC. A: relative PC is preferable because saved execution state remains relocation-independent; hardware translates on fetch.}
\stemrow[分页位数计算]{P7.12 Paging bits}{Stem: \(2^{32}\) physical bytes, page \(2^{10}\), \(2^{16}\) logical pages. A: logical address 26 bits; frame 1024 bytes; frame field 22 bits; table \(2^{16}\) entries; PTE 23 bits with valid bit.}
\stemrow[分页/分段翻译]{P7.13 Translate}{Stem: logical \code{0001010010111010}. A: paging: page 20, offset 186, frame 5 gives \code{0000010110111010}. Segmentation: seg 5, offset 186, base \(4+4096\times5\), physical \(20670\) = \code{0101000010111110}.}
\stemrow[段表越界]{P7.14 Seg table}{Stem: translate \((seg,offset)\). A: (0,198)\(\to858\); (2,156)\(\to378\); (1,530)\(\to\)fault; (3,444)\(\to1440\); (0,222)\(\to882\).}
\stemrow[紧凑收益阈值]{P7.15 Compaction}{Stem: compaction fraction. A: \(F\ge(1-f)/(1+kf)\), \(k=t/(2s)-1\). For \(f=.2,t=1000,s=50\), \(k=9\), so \(F\ge0.8/2.8\approx0.286\).}
\end{stemlist}

\cheatsection{Chapter 7 Fast Recall}
\begin{itemize}
  \item Fixed partitions waste inside; dynamic partitions waste between holes.
  \item Compaction fights external fragmentation but costs CPU and needs relocation.
  \item Buddy allocation rounds up to a power of two, splits downward, and coalesces upward.
  \item Paging maps page number to frame number; the offset is copied unchanged.
  \item Simple paging requires all pages resident; virtual-memory paging can load pages on demand.
\end{itemize}

\cheatsection{8 Virtual Memory}
\cheatsubsection{Core Idea and Locality}
\term[虚拟内存/只装一部分]{Virtual memory}: lets a process use a large logical address space while only part of it is in real/main memory; the rest is on secondary storage.\\
\term[虚拟地址/CPU生成]{Logical/virtual address}: generated by the CPU. \term[实地址/主存位置]{Physical/real address}: actual main-memory location. \term[地址转换]{Address translation}: MMU/OS mapping from virtual to physical.\\
\term[驻留集/在内存的页]{Resident set}: pages of a process currently in memory. \term[缺页/页不在内存]{Page fault}: referenced page is not resident, so the OS must fetch it.\\
\term[局部性]{Locality}: references cluster in time and space; virtual memory works because the current locality can fit in a modest resident set. \term[抖动/忙于换页]{Thrashing}: page faults dominate execution.

\cheatsubsection{Paging, TLB, Faults}
\[
\begin{aligned}
VA&=(p,d),\quad PA=(f,d)\\
p&=\lfloor VA/s\rfloor,\quad d=VA\bmod s\\
PA&=f\times s+d
\end{aligned}
\]
Offset \(d\) is copied unchanged; the page table maps page \(p\) to frame \(f\).\\
\term[页表项字段]{PTE fields}: present/valid bit, modify/dirty bit, reference/use bit, frame number, and protection/control bits.\\
\term[快表/缓存PTE]{TLB}: small fast cache of recent PTEs. \term[快表未命中]{TLB miss} means the PTE is not cached; \term[缺页中断]{page fault} means the page is not in memory.\\
\term[缺页处理流程]{Fault handling}: trap to OS, block process, read page from disk, replace a victim if needed, update PTE/TLB, ready the process, restart the faulting instruction.

\cheatsubsection{Policies}
\term[取页策略]{Fetch}: demand paging loads only on reference; prepaging also loads predicted nearby/future pages.\\
\term[放置策略]{Placement}: choose a free frame; with paging, placement is usually simple because any frame works.\\
\term[置换策略]{Replacement}: choose the page least likely to be used soon.
\begin{tabularx}{\linewidth}{@{}TY@{}}
\toprule
\term[最优/未来最远]{OPT} & evict page used farthest in future; optimal benchmark, not implementable.\\
\term[最近最少使用]{LRU} & evict least recently used page; good locality match, costly exact implementation.\\
\term[先进先出]{FIFO} & evict oldest resident page; simple, can evict hot pages and has Belady anomaly.\\
\term[时钟/二次机会]{Clock} & FIFO plus use bit; use \(=1\) gets a second chance, use \(=0\) is victim.\\
\term[增强时钟/优先干净页]{Enhanced clock} & prefer \((u,m)=(0,0)\), then \((0,1)\), then recently used pages; clean pages are cheaper to evict.\\
\bottomrule
\end{tabularx}
\term[驻留集管理/给几帧]{Resident-set management}: decides how many frames and local/global scope. \term[工作集/近期访问页]{Working set}: pages referenced in recent window \(\Delta\); keep it resident to avoid thrashing.\\
\term[清除策略/写脏页]{Cleaning}: demand cleaning writes dirty pages only on eviction; precleaning writes dirty pages early, possibly in batches.\\
\term[页大小权衡]{Page size}: small pages reduce internal fragmentation but enlarge page tables/TLB pressure; large pages reduce table overhead but may waste memory and weaken locality.

\cheatsection{Chapter 8 Review Q\&A}
\begin{stemlist}
\stemrow[简单分页vs虚存]{R8.1 Simple vs VM paging}{Stem: difference. A: simple paging requires all pages resident; virtual-memory paging keeps only needed pages and uses page faults.}
\stemrow[抖动定义]{R8.2 Thrashing}{Stem: explain. A: excessive page faults make the system spend most time paging instead of executing.}
\stemrow[局部性意义]{R8.3 Locality}{Stem: why crucial. A: references cluster, so a small resident set can cover the active working set.}
\stemrow[PTE字段]{R8.4 PTE fields}{Stem: typical elements. A: present/valid, modify/dirty, reference/use, frame number, protection/control.}
\stemrow[TLB作用]{R8.5 TLB}{Stem: purpose. A: cache recent PTEs to avoid an extra memory access for every translation.}
\stemrow[请求分页/预取]{R8.6 Fetch policies}{Stem: alternatives. A: demand paging loads on fault; prepaging loads likely-needed pages in advance.}
\stemrow[驻留集vs置换]{R8.7 Resident/replacement}{Stem: difference. A: resident-set management chooses frame count/scope; replacement chooses the victim page.}
\stemrow[Clock=FIFO+引用位]{R8.8 FIFO/clock}{Stem: relationship. A: clock is FIFO with a reference bit and second chance.}
\stemrow[页缓冲目的]{R8.9 Page buffering}{Stem: purpose. A: keep replaced pages on free/modified lists for quick reuse and batched dirty writes.}
\stemrow[固定+全局矛盾]{R8.10 Fixed+global}{Stem: why impossible. A: global replacement changes another process's frame count, violating fixed allocation.}
\stemrow[驻留集vs工作集]{R8.11 Resident/working set}{Stem: difference. A: resident set is pages actually in memory; working set is pages used recently.}
\stemrow[按需清除vs预清除]{R8.12 Cleaning}{Stem: demand cleaning vs precleaning. A: write dirty page at eviction vs write dirty pages before they are selected.}
\end{stemlist}

\cheatsection{Chapter 8 Problem Q\&A}
\begin{stemlist}
\stemrow[虚拟到物理]{P8.1 Translate}{Stem: page size 1024, given PTEs; VA 1052, 2221, 5499. A: use \(p=\lfloor a/1024\rfloor,d=a\bmod1024\); results \(7196\), page fault, \(379\).}
\stemrow[数组局部性]{P8.2 Arrays}{Stem: \(64\times64\) arrays, 1 KB pages, 4-page working set. A: column order gives 3 data faults per 4 assignments, total 3072; swap loops to row order gives 3 per 256, total 48.}
\stemrow[页表空间]{P8.3 Table space}{Stem: Fig. 8.3 page table and hashed inverted table. A: normal user page table \(=2^{20}\times4=4\) MB; hashed table \(=64\times4\times46\) bits \(=1472\) bytes, about 1.5 KB.}
\stemrow[置换算法计数]{P8.4 Replacement}{Stem: refs \(7,0,1,2,0,3,0,4,2,3,0,3,2\), 3 frames, count after init. A: FIFO 7 faults/70\%; LRU 6/60\%; Clock 6/60\%; OPT 4/40\%.}
\stemrow[Belady异常]{P8.5 FIFO anomaly}{Stem: refs A B C D A B E A B C D E. A: FIFO transfers: 3 frames \(=9\), 4 frames \(=10\); this is Belady anomaly.}
\stemrow[LRU/FIFO同分]{P8.6 LRU vs FIFO}{Stem: 33-reference trace, 4 frames. A: both LRU and FIFO have 17 faults, 16 hits, hit ratio \(16/33\approx48.5\%\); equality is trace-specific.}
\stemrow[页表也可缺页]{P8.7 VAX tables}{Stem: user page tables in virtual memory. A: saves real memory; disadvantage is possible page fault while accessing the page table itself.}
\stemrow[多级页表层数]{P8.10 Levels}{Stem: 64-bit address, 4 KB pages, 4-byte PTEs. A: offset 12 bits, page number 52 bits, 10 index bits per table page, so \(\lceil52/10\rceil=6\) levels.}
\stemrow[有效访问时间]{P8.11 EMAT}{Stem: 200 ns memory, 20 ns MMU overhead, 85\% TLB hit. A: no TLB \(=400\) ns; with TLB \(=.85(220)+.15(420)=250\) ns; \(EMAT=420-200h\).}
\stemrow[缺页上下界]{P8.12 Bounds}{Stem: length \(P\), \(N\) distinct pages, \(M\) frames. A: lower bound \(N\); upper bound \(P\), or \(N\) if \(M\ge N\).}
\stemrow[时钟算法类比]{P8.13 Snowplow}{Stem: circular snowplow analogy. A: clock algorithm; pointer clears reference bits as it moves, evicting pages not recently refreshed.}
\stemrow[引用位近似LRU]{P8.14 Ref-bit LRU}{Stem: use only reference bit. A: approximate LRU with aging: periodically sample/clear reference bits and evict the lowest age counter.}
\stemrow[地址格式位数]{P8.18 Address format}{Stem: 32 pages of 2 KB into 1 MB physical memory. A: logical address \(=5\)-bit page + \(11\)-bit offset; table \(=32\) entries \(\times9\)-bit frame; half memory makes frame field 8 bits, length unchanged.}
\end{stemlist}

\cheatsection{Chapter 8 Fast Recall}
\begin{itemize}
  \item Virtual memory = paging plus run-time translation plus page-fault recovery.
  \item Locality keeps the working set small; poor resident-set sizing causes thrashing.
  \item TLB miss is a cache miss; page fault is a residency miss.
  \item Replacement quality: OPT benchmark, LRU strong but costly, FIFO simple but weak, Clock practical.
  \item Clean victims avoid write-back I/O; dirty victims must be written first.
\end{itemize}

\cheatsection{9 Uniprocessor Scheduling}
\cheatsubsection{Schedulers and Criteria}
\term[长程/准入作业]{Long-term scheduler}: admits jobs to the system and controls the degree of multiprogramming.\\
\term[中程/交换挂起]{Medium-term scheduler}: swaps processes in/out of main memory and manages suspended states.\\
\term[短程/派遣CPU]{Short-term scheduler / dispatcher}: chooses the next ready process to run; invoked most often.\\
\term[调度评价指标]{Criteria}: response time, turnaround time, throughput, processor utilization, fairness, predictability, deadlines.
\[
T_r=finish-arrival,\qquad NTT=\frac{T_r}{T_s}
\]
\term[优先级/数值小级高]{Priority}: in this text, a smaller priority value usually means higher priority. \term[老化/防饥饿]{Aging} raises long-waiting jobs to prevent starvation.

\cheatsubsection{Decision Modes and Algorithms}
\term[非抢占]{Nonpreemptive}: current process runs until completion or blocking. \term[抢占]{Preemptive}: OS can interrupt the running process.\\
\begin{stemlist}
\term[先来先服务]{FCFS/FIFO} -- \textbf{Full name}: First-Come-First-Served / First-In-First-Out.\\
\textbf{Idea}: nonpreemptive; run the ready process that arrived first or waited longest.\\
\textbf{Pros}: simplest policy; little overhead; fair by arrival order.\\
\textbf{Cons}: convoy effect; short or I/O-bound jobs can wait behind one long CPU-bound job.

\term[轮转/时间片]{RR} -- \textbf{Full name}: Round Robin.\\
\textbf{Idea}: preemptive; each ready process gets a fixed time quantum \(q\), then goes to the queue tail if unfinished.\\
\textbf{Pros}: good response time and fairness for time-sharing systems.\\
\textbf{Cons}: small \(q\) causes many context switches; large \(q\) degenerates toward FCFS.

\term[最短作业优先]{SPN/SJF} -- \textbf{Full name}: Shortest Process Next / Shortest Job First.\\
\textbf{Idea}: nonpreemptive; choose the process with the smallest estimated service time.\\
\textbf{Pros}: often minimizes average waiting time when jobs arrive together; favors short jobs.\\
\textbf{Cons}: needs burst-time estimates; long jobs may starve.

\term[最短剩余时间]{SRT} -- \textbf{Full name}: Shortest Remaining Time.\\
\textbf{Idea}: preemptive version of SPN; always run the process with the least remaining CPU time.\\
\textbf{Pros}: very good response for short/new short jobs; low average turnaround in many workloads.\\
\textbf{Cons}: needs remaining-time estimates; more preemption overhead; long jobs may starve.

\term[最高响应比]{HRRN} -- \textbf{Full name}: Highest Response Ratio Next.\\
\textbf{Idea}: nonpreemptive; choose largest \(R=(w+s)/s=1+w/s\), where \(w\)=waiting time and \(s\)=service time.\\
\textbf{Pros}: short jobs score well, and long-waiting jobs age upward; reduces starvation compared with SPN.\\
\textbf{Cons}: must recompute ratios; still needs service-time estimates.

\term[反馈/多级反馈队列]{Feedback} -- \textbf{Full name}: Feedback / Multilevel Feedback Queue.\\
\textbf{Idea}: preemptive; new jobs start in high-priority queues, and jobs that use full quanta are demoted.\\
\textbf{Pros}: no service-time estimate needed; favors short, interactive, and I/O-bound jobs.\\
\textbf{Cons}: CPU-bound jobs can be pushed down and starve unless aging or longer low-level quanta are used.
\end{stemlist}

\cheatsection{Chapter 9 Q\&A Symbols}
\begingroup
\cheattinyfont
\renewcommand{\arraystretch}{0.74}
\begin{tabularx}{\linewidth}{@{}>{\raggedright\arraybackslash}p{0.16\linewidth}>{\raggedright\arraybackslash}p{0.25\linewidth}>{\raggedright\arraybackslash}p{0.17\linewidth}Y@{}}
\toprule
\term{缩写} & \term{English} & \term{中文} & \term{含义 / 公式}\\
\midrule
\term{AT} & Arrival Time & 到达时间 & 进入 ready queue 的时间\\
\term{BT / ST} & Burst Time / Service Time & 执行/服务时间 & 需要 CPU 运行的总时间\\
\term{Start Time} & Start Time & 开始时间 & 第一次获得 CPU 的时间\\
\term{CT / FT} & Completion / Finish Time & 完成时间 & 执行完毕的时间\\
\term{TAT} & Turnaround Time & 周转时间 & \code{CT - AT}\\
\term{WT} & Waiting Time & 等待时间 & \code{TAT - BT}\\
\term{RT} & Response Time & 响应时间 & \code{Start Time - AT}\\
\term{NTAT} & Normalized Turnaround Time & 归一化周转时间 & \code{TAT / BT}\\
\bottomrule
\end{tabularx}
\endgroup

\cheatsection{Chapter 9 Review Q\&A}
\begin{stemlist}
\stemrow[长中短调度]{R9.1 Scheduler types}{Stem: three types. A: long-term admits jobs, medium-term swaps processes, short-term dispatches a ready process to CPU.}
\stemrow[交互看响应]{R9.2 Interactive OS}{Stem: critical performance requirement. A: response time.}
\stemrow[周转vs响应]{R9.3 Turnaround/response}{Stem: difference. A: turnaround is submit/arrival to finish; response is submit/request to first visible response.}
\stemrow[数值小优先级高]{R9.4 Priority value}{Stem: low value means what. A: in this text/UNIX style, lower numeric value means higher priority.}
\stemrow[抢占vs非抢占]{R9.5 Preemption}{Stem: preemptive vs nonpreemptive. A: preemptive can interrupt a running process; nonpreemptive waits for finish or block.}
\stemrow[先来先服务]{R9.6 FCFS}{Stem: define. A: ready process waiting longest is selected first.}
\stemrow[时间片轮转]{R9.7 RR}{Stem: define. A: each ready process gets one time quantum, then returns to the tail if unfinished.}
\stemrow[最短进程]{R9.8 SPN}{Stem: define. A: nonpreemptively select the process with shortest expected service time.}
\stemrow[最短剩余]{R9.9 SRT}{Stem: define. A: preemptive SPN; select the process with shortest remaining time.}
\stemrow[响应比公式]{R9.10 HRRN}{Stem: define. A: select largest \((w+s)/s\), balancing short jobs with aging.}
\stemrow[多级反馈]{R9.11 Feedback}{Stem: define. A: multilevel feedback queues demote jobs that consume CPU quanta and favor short/I/O-bound jobs.}
\end{stemlist}

\cheatsection{Chapter 9 Problem Q\&A}
\begin{stemlist}
\par\noindent\term[甘特图+平均等待]{P9.1 SRT/priority/RR}: \textbf{Question}: workload below. Show SRT, nonpreemptive priority(smaller number = higher priority), RR \(q=30\) ms; compute average WT.
\begin{tabularx}{\linewidth}{@{}lrrrr@{}}
\toprule
P & P1 & P2 & P3 & P4\\
\midrule
BT & 50 & 20 & 100 & 40\\
Priority & 4 & 1 & 3 & 2\\
AT & 0 & 20 & 40 & 60\\
\bottomrule
\end{tabularx}
\textbf{Formula}: \(WT=FT-AT-BT\).\\
\textbf{SRT}: compare remaining time at arrivals. Gantt: \(0\!-\!20:P1\,|\,20\!-\!40:P2\,|\,40\!-\!70:P1\,|\,70\!-\!110:P4\,|\,110\!-\!210:P3\). FT \(=(70,40,210,110)\); WT \(=(20,0,70,10)\); avg \(=25\) ms.\\
\textbf{Priority}: nonpreemptive, so P1 finishes first; then priority order P2, P4, P3. Gantt: \(0\!-\!50:P1\,|\,50\!-\!70:P2\,|\,70\!-\!110:P4\,|\,110\!-\!210:P3\). WT \(=(0,30,70,10)\); avg \(=27.5\) ms.\\
\textbf{RR \(q=30\)}: unfinished process returns to queue tail. Gantt: \(0\!-\!30:P1\,|\,30\!-\!50:P2\,|\,50\!-\!70:P1\,|\,70\!-\!100:P3\,|\,100\!-\!130:P4\,|\,130\!-\!160:P3\,|\,160\!-\!170:P4\,|\,170\!-\!210:P3\). WT \(=(20,10,70,70)\); avg \(=42.5\) ms.

\par\noindent\term[多算法对比]{P9.2 Algorithm comparison}: \textbf{Question}: processes \(A(AT=0,BT=3), B(1,5), C(3,2), D(9,5), E(12,5)\). Compute Gantt, finish time, mean TAT, mean NTAT.
\textbf{Formulas}: \(TAT=FT-AT,\quad NTAT=TAT/BT\).\\
\textbf{FCFS}: \(0\!-\!3:A\,|\,3\!-\!8:B\,|\,8\!-\!10:C\,|\,10\!-\!15:D\,|\,15\!-\!20:E\). FT \(=(3,8,10,15,20)\); TAT \(=(3,7,7,6,8)\); mean \(=6.2\); mean NTAT \(=1.74\).\\
\textbf{RR \(q=1\)}: unit Gantt: A B A B C A B C B D / B D E D E D E D E E. FT \(=(6,11,8,18,20)\); TAT \(=(6,10,5,9,8)\); mean \(=7.6\); mean NTAT \(=1.98\).\\
\textbf{RR \(q=4\)}: \(0\!-\!3:A\,|\,3\!-\!7:B\,|\,7\!-\!9:C\,|\,9\!-\!10:B\,|\,10\!-\!14:D\,|\,14\!-\!18:E\,|\,18\!-\!19:D\,|\,19\!-\!20:E\). FT \(=(3,10,9,19,20)\); mean TAT \(=7.2\); mean NTAT \(=1.88\).\\
\textbf{SPN/SRT}: \(0\!-\!3:A\,|\,3\!-\!5:C\,|\,5\!-\!10:B\,|\,10\!-\!15:D\,|\,15\!-\!20:E\). FT \(=(3,10,5,15,20)\); mean TAT \(=5.6\); mean NTAT \(=1.32\). SRT same here because no later job has shorter remaining time.\\
\textbf{HRRN}: choose largest \((w+s)/s\). Gantt: \(0\!-\!3:A\,|\,3\!-\!8:B\,|\,8\!-\!10:C\,|\,10\!-\!15:D\,|\,15\!-\!20:E\). FT \(=(3,8,10,15,20)\); mean TAT \(=6.2\); mean NTAT \(=1.74\). Feedback results depend on exact queue/demotion rule; one strict \(q=1\) convention gives FT \(=(7,18,6,19,20)\), mean TAT \(=9.0\), mean NTAT \(=2.17\).
\stemrow[SPN最优交换证明]{P9.3 SPN proof}{Stem: prove minimum average wait when all jobs arrive together. A: exchange adjacent jobs \(i,j\); if \(s_i\le s_j\), order \(i,j\) contributes \(2T+s_i\) wait, while \(j,i\) gives \(2T+s_j\). Repeated swaps yield shortest-first.}
\stemrow[指数平均预测]{P9.4 Exponential average}{Stem: bursts \(6,4,6,4,13,13,13\), initial 10. A: \(S_{n+1}=\alpha T_n+(1-\alpha)S_n\); larger \(\alpha\) follows recent bursts faster. For \(\alpha=.5\): \(8,6,6,5,9,11,12\); for \(\alpha=.8\): \(6.8,4.56,5.712,4.342,11.268,12.654,12.931\).}
\stemrow[估计参数]{P9.5 Bounded estimate}{Stem: roles of \(\alpha,\beta\) in bounded SPN estimate. A: \(\alpha\) controls smoothing vs responsiveness; \(\beta\) scales the estimate before clamping to lower/upper bounds, biasing a job to look shorter or longer.}
\stemrow[反馈抢占细节]{P9.6 Feedback variant}{Stem: why A may run 2 or 3 units before B. A: if a new high-priority job immediately preempts a lower-priority running job, A stops after 2; if rescheduling waits for the current quantum to expire, A may run 3.}
\stemrow[最小最大响应比]{P9.7 Minimax RR}{Stem: derive minimax response-ratio scheduling. A: build the order backward; at each step, choose as last the job with the smallest response ratio at the current final completion time, remove it, and repeat.}
\stemrow[逆序归纳证明]{P9.8 Minimax proof}{Stem: prove preceding algorithm. A: any schedule's last job has ratio fixed by the final time; choosing the minimum possible last ratio gives the best upper bound, then induction applies to the remaining jobs.}
\stemrow[排队论周转]{P9.9 FIFO residence}{Stem: show FIFO \(T_r=T_s/(1-\rho)\). A: M/M/1 gives \(W_q=\rho T_s/(1-\rho)\); add service time, so \(T_r=T_s+W_q=T_s/(1-\rho)\).}
\stemrow[理想RR伸缩]{P9.10 Ideal RR}{Stem: processor-sharing RR with infinitesimal quantum. A: Poisson/exponential processor sharing gives conditional mean response \(R_x=x/(1-\rho)\); RR stretches service by factor \(1/(1-\rho)\).}
\stemrow[加权RR]{P9.11 Duplicate PCB pointers}{Stem: same process appears twice in RR queue. A: it gets two quanta per cycle; advantage is weighted RR; equivalent design is a weight field or enlarged effective quantum.}
\stemrow[Selfish RR公式]{P9.12 Selfish RR}{Stem: derive selfish RR response time. A: with \(\rho=\lambda s\), \(\rho'=\rho(1-b/a)\), \(0\le b<a\), \(R_x=\frac{s}{1-\rho}+\frac{x-s}{1-\rho'}\).}
\stemrow[动态时间片风险]{P9.13 Dynamic RR quantum}{Stem: divide max response time by ready-process count. A: only a rough heuristic; changing arrivals, swap time, and dispatch overhead can break the response guarantee and make quanta too small.}
\stemrow[反馈偏I/O型]{P9.14 Feedback favors}{Stem: CPU-bound or I/O-bound. A: I/O-bound/interactive processes, because they block before using full quanta and avoid demotion.}
\stemrow[优先级反转]{P9.15 Dekker danger}{Stem: static priority plus busy waiting. A: priority inversion/starvation: a high-priority spinner can keep the low-priority process from running and leaving the critical section.}
\stemrow[批处理周转时间]{P9.16 Batch turnaround}{Stem: jobs A--E times \(15,9,3,6,12\), priorities \(6,3,7,9,4\). A: RR \(q=1\): A45 B35 C13 D26 E42 avg 32.2; priority order B,E,A,C,D gives A36 B9 C39 D45 E21 avg 30; FCFS gives A15 B24 C27 D33 E45 avg 28.8; SJF order C,D,B,E,A gives A45 B18 C3 D9 E30 avg 21.}
\end{stemlist}

\cheatsection{Chapter 9 Fast Recall}
\begin{itemize}
  \item Long-term controls admission; medium-term controls swapping; short-term chooses the next ready process.
  \item Turnaround \(=finish-arrival\); normalized turnaround \(=T_r/T_s\); waiting \(=finish-arrival-service\).
  \item RR needs a useful quantum: small improves response but raises overhead; large approaches FCFS.
  \item SPN/SRT improve average wait for short jobs but need burst estimates and may starve long jobs.
  \item HRRN adds aging through \(1+w/s\); feedback favors short and I/O-bound processes without knowing burst length.
\end{itemize}

\cheatsection{10 Multiprocessor, Multicore, Real-Time}
\cheatsubsection{Multiprocessor Scheduling}
\term[紧耦合多处理器]{Tightly coupled multiprocessor}: processors share common main memory and are controlled by one OS.\\
\term[粒度/同步频率]{Granularity}: synchronization frequency. Fine = instruction-level parallelism; medium = cooperating threads in one application; coarse = concurrent processes; very coarse = distributed nodes; independent = unrelated tasks with no explicit synchronization.\\
\term[处理器分配]{Assignment}: static assignment keeps a process on one processor, with low overhead but possible imbalance; dynamic assignment uses a common queue, improving balance but hurting cache locality.\\
\term[内核组织]{Kernel organization}: master/slave is simple but can bottleneck or fail at the master; peer scheduling scales better but needs stronger synchronization.\\
\term[多处理器派遣]{Dispatching}: with multiple processors, simple policies such as FCFS may be adequate because one long job need not block all other work.

\cheatsubsection{Thread and Core Scheduling}
\begingroup
\cheattinyfont
\begin{tabularx}{\linewidth}{@{}TY@{}}
\toprule
\term[技术]{Technique} & \term[思路/优缺点]{Idea / Pros / Cons}\\
\midrule
\term[负载共享]{Load sharing} & Idea: global queue. Pros: simple, balanced, keeps CPUs busy. Cons: queue contention, cache-affinity loss, related threads separated.\\
\term[组调度/相关线程同跑]{Gang scheduling} & Idea: co-schedule related threads. Pros: less barrier/lock wait, lower communication delay. Cons: needs CPUs, may waste slots, harder with mixed jobs.\\
\term[专用分配]{Dedicated assignment} & Idea: reserve processors per app. Pros: strong affinity, low dispatch cost, predictable progress. Cons: idle when threads block, poor sharing, bad if thread count varies.\\
\term[动态调度]{Dynamic scheduling} & Idea: vary thread count and processors at run time. Pros: adapts to load, high utilization, can favor urgent work. Cons: complex policy, migration overhead, less predictable.\\
\term[核心放置/亲和性]{Core placement} & Idea: choose cores by cache/memory behavior. Pros: shared-data threads share cache; competitors separated. Cons: needs workload knowledge; may hurt load balance.\\
\bottomrule
\end{tabularx}
\endgroup
\term[多核经验法则]{Multicore rule}: co-locate data-sharing threads; separate threads that compete for cache or memory bandwidth.

\cheatsubsection{Real-Time Basics}
\term[硬实时/错过不可接受]{Hard real-time}: missing a deadline can cause unacceptable damage or fatal failure.\\
\term[软实时/迟到仍可用]{Soft real-time}: deadlines are desirable but not mandatory; late results may still have value.\\
\term[周期任务]{Periodic task}: occurs repeatedly at regular intervals. \term[非周期任务]{Aperiodic task}: has no fixed interval but has a start or finish deadline.\\
\term[实时OS要求]{RTOS requirements}: determinism, responsiveness, user control, reliability, and fail-soft operation.

\cheatsection{Chapter 10 Review Q\&A}
\begin{stemlist}
\stemrow[同步粒度]{R10.1 Granularity}{Stem: five synchronization granularities. A: fine, medium, coarse, very coarse, and independent parallelism; finer granularity means more frequent synchronization.}
\stemrow[四种线程调度]{R10.2 Thread scheduling}{Stem: four thread-scheduling techniques. A: load sharing, gang scheduling, dedicated processor assignment, and dynamic scheduling.}
\stemrow[负载共享变体]{R10.3 Load sharing versions}{Stem: three load-sharing versions. A: FCFS shared queue; smallest number of threads first; preemptive smallest number of threads first.}
\stemrow[硬实时vs软实时]{R10.4 Hard vs soft RT}{Stem: hard vs soft real-time tasks. A: hard tasks must meet deadlines; soft tasks prefer deadlines but late completion can still be useful.}
\stemrow[周期vs非周期]{R10.5 Periodic vs aperiodic}{Stem: periodic vs aperiodic tasks. A: periodic tasks recur every period \(T\); aperiodic tasks arrive irregularly but still have timing constraints.}
\stemrow[RTOS五要求]{R10.6 RTOS requirements}{Stem: five RTOS requirement areas. A: deterministic behavior, fast response, user control over priorities/resources, high reliability, and graceful fail-soft degradation.}
\stemrow[实时调度类别]{R10.7 RT algorithms}{Stem: four real-time scheduling classes. A: static table-driven, static priority-driven preemptive, dynamic planning-based, and dynamic best-effort scheduling.}
\stemrow[实时任务参数]{R10.8 Task data}{Stem: task information useful for real-time scheduling. A: ready time, starting/completion deadline, processing time, resource needs, priority, periodicity, and subtask/dependency structure.}
\end{stemlist}

\cheatsection{Chapter 10 Fast Recall}
\begin{itemize}
  \item Static assignment favors affinity; dynamic assignment favors load balance.
  \item Load sharing is simple but can lose cache affinity and separate cooperating threads.
  \item Gang scheduling is best when related threads synchronize often.
  \item Dedicated assignment reduces switching but may waste processors when threads block.
  \item Multicore scheduling balances load, shared-cache benefit, and resource contention.
\end{itemize}

\cheatsection{12 File Management}
\cheatsubsection{Core Model and Organizations}
\term[文件系统]{File system}: maps named files and logical records onto secondary-storage blocks. Logical hierarchy: field \(\to\) record \(\to\) file \(\to\) database. Users see files/records; devices transfer blocks.\\
\term[文件管理系统]{File management system}: OS software for create, delete, open, close, read, write, protect, share, back up, and recover files.\\
\term[文件组织/逻辑访问]{File organization}: logical record layout/access method; choose by access time, update ease, storage economy, maintenance simplicity, and reliability.
\begingroup
\cheattinyfont
\renewcommand{\arraystretch}{0.68}
\begin{tabularx}{\linewidth}{@{}>{\raggedright\arraybackslash}p{0.31\linewidth}Y@{}}
\term[访问方法层]{Access method} & top; sequential, indexed, hashed access.\\
\term[逻辑I/O/记录级]{Logical I/O} & record-level operations.\\
\term[基本I/O管理]{Basic I/O supervisor} & I/O requests, buffers, scheduling.\\
\term[物理块I/O]{Basic file system / Physical I/O} & block-level physical I/O.\\
\term[设备驱动]{Device drivers} & lowest; control storage devices.\\
\end{tabularx}
\endgroup
\begin{tabularx}{\linewidth}{@{}TY@{}}
\toprule
\term[堆文件/追加无序]{Pile} & unordered arrival order; simple append, slow search.\\
\term[顺序文件/按键排序]{Sequential} & records sorted by key; strong for batch/full scans.\\
\term[索引顺序文件]{Indexed sequential} & sorted file plus index and overflow; faster lookup and sequential processing.\\
\term[索引文件/多索引]{Indexed} & one or more indexes; flexible random access, index overhead.\\
\term[直接/散列文件]{Direct/hashed} & hash key to location; fastest exact-key lookup, weaker for ordered scans.\\
\bottomrule
\end{tabularx}

\cheatsubsection{Directories, B-Trees, Blocking}
\term[目录/文件元数据表]{Directory}: a file containing file metadata: name/type/org, volume/start/size, owner/access rights, timestamps, open/lock state. Ops: search, create, delete, list, update.\\
\term[路径名/绝对相对]{Pathname}: absolute starts at root; relative is interpreted from the working/current directory. Tree directories avoid name conflicts by unique pathnames.\\
\term[访问权限]{Access rights}: none, knowledge, execution, reading, appending, updating, changing protection, deletion.\\
\term[B树/平衡多路索引]{B-tree}: balanced multiway search tree; all leaves same depth, broad and shallow. Minimum degree \(d\): max keys \(2d-1\), max children \(2d\), nonroot min keys \(d-1\). On full-node insert, split and promote median.
\begingroup
\cheattinyfont
\renewcommand{\arraystretch}{0.74}
\begin{tabularx}{\linewidth}{@{}>{\raggedright\arraybackslash}p{0.34\linewidth}Y@{}}
\toprule
\term[固定分块]{Fixed blocking} & Idea: fixed records, integral records/block. Pros: simplest, easy addressing. Cons: internal fragmentation; poor for variable records.\\
\term[可变跨块]{Variable spanned} & Idea: variable records may cross blocks. Pros: best space use; record may exceed block. Cons: complex I/O/update; one record may need two block reads.\\
\term[可变不跨块]{Variable unspanned} & Idea: variable records cannot cross blocks. Pros: simpler; one record stays in one block. Cons: leftover waste; record \(\le\) block.\\
\bottomrule
\end{tabularx}
\endgroup

\cheatsubsection{Allocation and Free Space}
\begingroup
\cheattinyfont
\renewcommand{\arraystretch}{0.74}
\begin{tabularx}{\linewidth}{@{}>{\raggedright\arraybackslash}p{0.31\linewidth}Y@{}}
\toprule
\term[连续分配]{Contiguous} & Idea: start block + length. Pros: fastest sequential/direct access; tiny metadata. Cons: external fragmentation; hard growth; size often needed early.\\
\term[链接分配]{Chained} & Idea: each block points to next. Pros: no external fragmentation; easy growth. Cons: slow random access; pointer overhead; poor locality.\\
\term[索引分配]{Indexed} & Idea: index block lists data blocks. Pros: good random/sequential access; flexible growth. Cons: index overhead; small-file waste; large files need multilevel indexes.\\
\bottomrule
\end{tabularx}
\term[预分配]{Preallocation}: reserve max size early. Pros: simple, may be contiguous/fast. Cons: wastes space; max size hard to know.\\
\term[动态分配]{Dynamic allocation}: allocate as file grows. Pros: space efficient, flexible. Cons: more metadata; may fragment.\\
\term[首次适配]{First fit}: first large-enough free portion. Pros: fast. Cons: small holes near front. \term[最佳适配]{Best fit}: smallest large-enough portion. Pros: least immediate leftover. Cons: slow search, tiny fragments. \term[邻近适配]{Nearest fit}: closest to prior file allocation. Pros: locality. Cons: not space-optimal; needs position state.\\
\endgroup
\term[空闲空间管理]{Free-space management}: bit table/bitmap, chained free portions, indexing, free block list. Bitmap bytes \(=disk\ bytes/(8\times block\ size)\).\\
\term[易混点]{Do not confuse}: indexed file = logical file organization; indexed allocation = physical block allocation.

\cheatsection{Chapter 12 Review Q\&A}
\begin{stemlist}
\stemrow[字段vs记录]{R12.1 Field/record}{Stem: difference. A: a field is one data item; a record is related fields treated as one unit.}
\stemrow[文件vs数据库]{R12.2 File/database}{Stem: difference. A: a file is similar records; a database is related data, often several files with explicit relationships.}
\stemrow[文件管理系统]{R12.3 FMS}{Stem: define file management system. A: OS software that provides file creation, deletion, access, modification, protection, sharing, backup, and recovery.}
\stemrow[组织选择标准]{R12.4 Organization criteria}{Stem: choosing file organization. A: short access time, ease of update, storage economy, simple maintenance, and reliability.}
\stemrow[五种文件组织]{R12.5 Five organizations}{Stem: list and define. A: pile = arrival/full scan; sequential = key-sorted/batch; indexed sequential = index+overflow/faster lookup; indexed = multiple indexes/fields; direct/hashed = hash key to location.}
\stemrow[索引顺序查找]{R12.6 Indexed sequential search}{Stem: why faster than sequential. A: the index jumps near the target, so only a small local sequential search remains.}
\stemrow[目录操作]{R12.7 Directory ops}{Stem: typical operations. A: search, create file, delete file, list directory, and update directory.}
\stemrow[路径和工作目录]{R12.8 Path/working dir}{Stem: relationship. A: a relative pathname is resolved from the working directory; an absolute pathname starts at root.}
\stemrow[访问权限列表]{R12.9 Access rights}{Stem: typical rights. A: none, knowledge, execution, reading, appending, updating, changing protection, deletion.}
\stemrow[三种记录分块]{R12.10 Blocking}{Stem: three methods. A: fixed; variable-length spanned; variable-length unspanned.}
\stemrow[三种块分配]{R12.11 Allocation}{Stem: three methods. A: contiguous, chained, and indexed file allocation.}
\end{stemlist}

\cheatsection{Chapter 12 Problem Q\&A}
\begin{stemlist}
\stemrow[分块因子公式]{P12.1 Blocking factor}{Stem: \(B,R,P\)=block/record/pointer sizes; \(F\)=blocking factor. A: fixed/unspanned \(F\approx\lfloor B/R\rfloor\); spanned \(F\approx(B-P)/R\).}
\stemrow[增长区段/FAT项]{P12.2 Growing portions}{Stem: portions \(1,2,4,\ldots\), \(n\) records, blocking factor \(F\). A: let \(m=\lceil n/F\rceil\) blocks; FAT entries \(\le\lceil\log_2(m+1)\rceil\); just after a new portion, unused space is at most that portion minus one block/record unit, less than half allocated.}
\stemrow[按场景选组织]{P12.3 Choose organization}{Stem: choose org for speed/storage/update: (a) infrequent update + frequent random access; (b) frequent update + frequent whole-file access; (c) frequent update + frequent random access. A: (a) direct/hashed, or indexed for multiple keys; (b) indexed sequential; (c) indexed, or hashed for one exact key.}
\stemrow[B树插入分裂]{P12.4 Insert 97}{Stem: B-tree insert. A: insert into full right leaf, split \([73,85,88,90,96,97]\) at 88; root also full, split at 51; new root \([51]\), right child includes \([61,71,88]\).}
\stemrow[自顶向下分裂]{P12.5 Top-down split}{Stem: split every full node while descending. A: advantage: no backtracking and simpler one-pass insertion; disadvantages: unnecessary splits, extra writes, lower occupancy, possible earlier height growth.}
\stemrow[B树高度上界]{P12.6 B-tree height}{Stem: upper bound on height. A: worst case has minimum occupancy, so \(n\ge2d^h-1\) and \(h\le\log_d((n+1)/2)\), with root-to-leaf edge height.}
\stemrow[最后块浪费]{P12.7 Last-block waste}{Stem: 16K blocks; \(S=41{,}600;640{,}000;4{,}064{,}000\). A: \(A=\lceil S/16K\rceil16K\), \(W\%=(A-S)/A\); \(15.36\%,2.34\%,0.383\%\).}
\stemrow[目录优点]{P12.8 Directory benefits}{Stem: advantages. A: symbolic names, name-to-location mapping, hierarchy, fewer name conflicts, access control, easier search/list/update.}
\stemrow[目录特殊/普通文件]{P12.9 Directory file type}{Stem: special file vs ordinary file. A: special protects format/integrity but is less flexible; ordinary gives uniform tools but risks corruption unless strongly protected.}
\stemrow[限制目录深度]{P12.10 Limited depth}{Stem: effect of shallow tree. A: users lose deep organization and long paths; OS gains simpler pathname parsing, bounded buffers/tables, and less recursive traversal.}
\stemrow[空闲指针丢失恢复]{P12.11 Lost free pointer}{Stem: hierarchical FS keeps free disk space in a free-space list: (a) if pointer to free space is lost, can system reconstruct the list? (b) scheme so pointer is never lost from one memory failure. A: scan directories/FAT/inodes, mark allocated blocks, rebuild unmarked free blocks; store redundant logged/checksummed pointer copies in stable storage.}
\stemrow[UNIX V最大文件]{P12.12 UNIX V max file}{Stem: In UNIX System V, block length is 1 KB and each block holds 256 block addresses; using the inode scheme, find the maximum file size. A: with 10 direct plus single/double/triple indirect, max \(=(10+256+256^2+256^3)\times1024=17{,}247{,}250{,}432\) bytes, about 16.06 GiB.}
\stemrow[inode容量/访问次数]{P12.13 8K inode}{Stem: UNIX inode: 12 direct pointers plus single/double/triple indirect; block=sector=8K; 32-bit disk pointer = 8-bit disk ID + 24-bit block ID. Find (a) max file size, (b) max partition size, (c) disk accesses for byte 13,423,956 when inode is already in memory. A: pointers/block \(=2048\); max file \(=(12+2048+2048^2+2048^3)8K=70{,}403{,}120{,}791{,}552\) bytes. Partition limit is 128 GiB per 24-bit block field per disk, or 32 TiB using all 8 disk bits. Byte 13,423,956 is in the single-indirect range, so read indirect block + data block = 2 disk accesses.}
\end{stemlist}

\end{multicols*}

\end{document}
